\documentclass[../main.tex]{subfiles}

\begin{document}
\chapter{Lie Algebras}

\section{Geometric Intuition}

We begin with the geometric origins of Lie algebras. The central idea is that a \textbf{Lie group describes continuous transformations}, while its \textbf{Lie algebra describes the infinitesimal behavior of these transformations}. The connection between the two is established by differentiation and the exponential map.

\subsection{Lie Groups and Their Tangent Spaces}

Recall that a Lie group $G$ is both a group and a smooth manifold, such that the group multiplication and inversion are smooth maps. For simplicity, we shall mainly consider matrix Lie groups.

The Lie algebra $\mathfrak{g}$ associated with $G$ can be identified with the tangent space at the identity:
\begin{equation}
  \mathfrak{g}=T_I G.
\end{equation}
More precisely, an element $X\in\mathfrak{g}$ can be regarded as the tangent vector of a smooth curve $A(t)\in G$ satisfying
\begin{equation}
  A(0)=I,\quad A'(0)=X.
\end{equation}
Geometrically, $X$ represents an infinitesimal transformation near the identity.

Conversely, given $X\in\mathfrak{g}$, the matrix exponential generates a one-parameter subgroup of $G$:
\begin{equation}
  A(t)=e^{tX}.
\end{equation}
Thus, the exponential map allows us to recover finite transformations from infinitesimal generators. In general, the exponential map need not cover the entire group, but it describes the group locally near the identity.

\begin{example}{$SO(2)$}{SO2 Geometry}
  Consider the rotation group
  \begin{equation}
    SO(2)=\left\{
      \begin{pmatrix}
        \cos\theta & -\sin\theta\\
        \sin\theta & \cos\theta
      \end{pmatrix}
      :\theta\in\mathbb{R}
    \right\}.
  \end{equation}
  Differentiating the rotation matrix at $\theta=0$, we obtain
  \begin{equation}
    X=\begin{pmatrix}
      0 & -1\\
      1 & 0
    \end{pmatrix}.
  \end{equation}
  The corresponding Lie algebra is
  \begin{equation}
    \mathfrak{so}(2)=\left\{
      \begin{pmatrix}
        0 & -a\\
        a & 0
      \end{pmatrix}
      :a\in\mathbb{R}
    \right\}.
  \end{equation}
  On the other hand, exponentiating the infinitesimal generator gives
  \begin{equation}
    e^{tX}=
    \begin{pmatrix}
      \cos t & -\sin t\\
      \sin t & \cos t
    \end{pmatrix}.
  \end{equation}
  Hence the elements of $SO(2)$ describe finite rotations, while the elements of $\mathfrak{so}(2)$ describe their infinitesimal generators.
\end{example}


In Figure~\ref{fig:rotations-and-tangent-line}, distinguish the identity $I$ in the group from $0$ in its Lie algebra: $\exp(0)=I$.

\begin{figure}[H]
    \centering
    \incfig{rotations-and-tangent-line}
    \caption{The circle parametrizes $SO(2)$ by its angle; it is not a circle of matrix entries. Its tangent direction at $I$ is $J$. The algebra is the line $\mathbb RJ$, and exponentiation wraps it around the group.}
    \label{fig:rotations-and-tangent-line}
\end{figure}

The Lie algebra is not merely a vector space. For matrix Lie groups, it is equipped with the Lie bracket
\begin{equation}
  [X,Y]=XY-YX.
\end{equation}
Geometrically, the Lie bracket measures the failure of two infinitesimal transformations to commute. For example, infinitesimal rotations around different axes in $SO(3)$ generally do not commute.

The commutator appears already in the first-order approximations:
\begin{equation*}
  (I + \epsilon X)(I + \delta Y) - (I + \delta Y)(I + \epsilon X) = \epsilon\delta(XY - YX)
\end{equation*}
Equivalently, expanding the group commutator near the identity gives
\begin{equation*}
  e^{\epsilon X} e^{\delta Y} e^{-\epsilon X} e^{-\delta Y} = I + \epsilon\delta(XY - YX) + O(\epsilon^2 \delta + \epsilon \delta^2).
\end{equation*}

\subsection{Homomorphisms and Differentiation}

Suppose $G$ and $H$ are Lie groups with corresponding Lie algebras $\mathfrak{g}$ and $\mathfrak{h}$. Let
\begin{equation}
  \Phi:G\to H
\end{equation}
be a Lie group homomorphism, i.e. a smooth map satisfying
\begin{equation}
  \Phi(AB)=\Phi(A)\Phi(B).
\end{equation}
Since $\Phi$ maps the identity of $G$ to the identity of $H$, its differential at the identity defines a linear map between the corresponding tangent spaces:
\begin{equation}
  \phi=(d\Phi)_I:\mathfrak{g}\to\mathfrak{h}.
\end{equation}
This is the associated Lie algebra homomorphism, which satisfies
\begin{equation}
  \phi([X,Y])=[\phi(X),\phi(Y)].
\end{equation}

An important relation between $\Phi$ and $\phi$ is
\begin{equation}
  \Phi(e^{tX})=e^{t\phi(X)}.
\end{equation}
Thus, applying a group homomorphism to a finite transformation generated by $X$ is equivalent to first applying the associated Lie algebra homomorphism to $X$, and then exponentiating.

\begin{example}{Angle Doubling}{Angle Doubling}
  Consider the Lie group homomorphism
  \begin{equation}
    \Phi:SO(2)\to SO(2),\quad
    R(\theta)\mapsto R(2\theta).
  \end{equation}
  The corresponding Lie algebra homomorphism is
  \begin{equation}
    \phi:\mathfrak{so}(2)\to\mathfrak{so}(2),
    \quad X\mapsto 2X.
  \end{equation}
  Hence the group homomorphism doubles the finite rotation angle, while its differential doubles the infinitesimal generator.
\end{example}

\begin{remark}
  The kernel of a Lie group homomorphism is a closed normal subgroup, and its Lie algebra is precisely the kernel of the associated Lie algebra homomorphism:
  \begin{equation}
    \operatorname{Lie}(\ker\Phi)=\ker\phi.
  \end{equation}
  Geometrically, $\ker\phi$ consists of all tangent directions at the identity along which $\Phi$ has zero first-order variation.
\end{remark}

\subsection{Conjugation and the Adjoint Representation}

For any $A\in G$, conjugation defines a group automorphism
\begin{equation}
  C_A:G\to G,\quad B\mapsto ABA^{-1}.
\end{equation}
Since $C_A(I)=I$, its differential at the identity maps the tangent space $\mathfrak{g}$ to itself. Explicitly, if $X=A'(0)$ for some smooth curve $A(t)$ through the identity, then
\begin{equation}
  (dC_B)_I(X)=BXB^{-1}
\end{equation}
for any fixed $B\in G$.

This motivates the following definition.

\begin{definition}{Adjoint Representation of Lie Groups}{Adjoint Representation of Lie Groups}
  Let $G$ be a matrix Lie group with Lie algebra $\mathfrak{g}$. For each $A\in G$, we define the \textbf{adjoint transformation}
  \begin{equation}
    \Ad_A:\mathfrak{g}\to\mathfrak{g},
    \quad X\mapsto AXA^{-1}.
  \end{equation}
  The map
  \begin{equation}
    \Ad:G\to GL(\mathfrak{g}),
    \quad A\mapsto\Ad_A
  \end{equation}
  is called the \textbf{adjoint representation} of $G$.
\end{definition}

Here $GL(\mathfrak{g})$ denotes the group of all invertible linear transformations on the vector space $\mathfrak{g}$. Notice that $\Ad_A$ is itself a linear transformation, whereas $\Ad$ assigns such a transformation to each element $A\in G$.

\begin{proposition}{}{}
  The adjoint representation $\Ad:G\to GL(\mathfrak{g})$ is a Lie group homomorphism.
\end{proposition}
\begin{proof}
  For any $A,B\in G$ and $X\in\mathfrak{g}$, we have
  \begin{equation}
    \begin{aligned}
      \Ad_{AB}(X)
      &=(AB)X(AB)^{-1}\\
      &=A(BXB^{-1})A^{-1}\\
      &=\Ad_A(\Ad_B(X)).
    \end{aligned}
  \end{equation}
  Therefore,
  \begin{equation}
    \Ad_{AB}=\Ad_A\circ\Ad_B.
  \end{equation}
  Smoothness follows from the smooth dependence of $AXA^{-1}$ on $A$.
\end{proof}

Geometrically, $\Ad_A$ describes how conjugation by $A$ transforms infinitesimal directions at the identity. In particular, for $SO(3)$, the adjoint action corresponds to rotating the axes of infinitesimal rotations.

\subsection{From $\Ad$ to $\ad$}

Since $\Ad:G\to GL(\mathfrak{g})$ is a Lie group homomorphism, we can differentiate it at the identity to obtain a Lie algebra homomorphism
\begin{equation}
  \ad:\mathfrak{g}\to\mathfrak{gl}(\mathfrak{g}),
\end{equation}
where $\mathfrak{gl}(\mathfrak{g})$ denotes the Lie algebra of all linear transformations on $\mathfrak{g}$, equipped with the commutator bracket.

For any $X,Y\in\mathfrak{g}$, we have
\begin{equation}
  \begin{aligned}
    \ad_X(Y)
    &=\left.\frac{d}{dt}\right|_{t=0}
      \Ad_{e^{tX}}(Y)\\
    &=\left.\frac{d}{dt}\right|_{t=0}
      e^{tX}Ye^{-tX}\\
    &=XY-YX\\
    &=[X,Y].
  \end{aligned}
\end{equation}

Thus, the Lie bracket naturally arises by differentiating the conjugation action of the Lie group.

\begin{remark}
  It is important to distinguish the two levels:
  \begin{equation}
    \begin{aligned}
      \Ad &:G\to GL(\mathfrak{g}),\\
      \ad &:\mathfrak{g}\to\mathfrak{gl}(\mathfrak{g}).
    \end{aligned}
  \end{equation}
  The first is a Lie group homomorphism, while the second is its associated Lie algebra homomorphism. Their actions are given by
  \begin{equation}
    \Ad_A(X)=AXA^{-1},\quad
    \ad_X(Y)=[X,Y].
  \end{equation}
  In particular, the identity
  \begin{equation}
    \ad_{[X,Y]}=[\ad_X,\ad_Y]
  \end{equation}
  is equivalent to the Jacobi identity.
\end{remark}

This overview connects tangent spaces, the exponential map, and differentiation of homomorphisms. We now introduce Lie algebras abstractly, then establish these relationships for matrix Lie groups.

\section{Definitions and Examples}

We now introduce the abstract notion of a Lie algebra:

\begin{definition}{Lie Algebras}{Lie Algebras}
  A \textbf{finite dimensional real or complex Lie algebra} is a finite dimensional vector space $\mathfrak{g}$ over $\mathbb{R}$ or $\mathbb{C}$ together with a map $[\cdot,\cdot]:\mathfrak{g}\times\mathfrak{g}\to\mathfrak{g}$ called the \textbf{Lie bracket} such that:
  \begin{itemize}
    \item $[\cdot,\cdot]$ is bilinear, i.e. for all $x,y,z\in\mathfrak{g}$ and $\alpha,\beta\in\mathbb{R}$ or $\mathbb{C}$, we have $[\alpha x+\beta y,z]=\alpha[x,z]+\beta[y,z]$ and $[z,\alpha x+\beta y]=\alpha[z,x]+\beta[z,y]$.
    \item $[\cdot,\cdot]$ is antisymmetric, i.e. for all $x,y\in\mathfrak{g}$, we have $[x,y]=-[y,x]$.
    \item $[\cdot,\cdot]$ satisfies the Jacobi identity, i.e. for all $x,y,z\in\mathfrak{g}$, we have
      \begin{equation}
        [x,[y,z]]+[y,[z,x]]+[z,[x,y]]=0.
      \end{equation}
  \end{itemize}
  We say two elements $x,y\in\mathfrak{g}$ \textbf{commute} if $[x,y]=0$. A Lie algebra is called \textbf{commutative} if all its elements commute.
\end{definition}

Antisymmetry implies that $[x,x]=0$ for all $x\in\mathfrak{g}$. The bracket need not be associative; the Jacobi identity plays the role of a compatibility condition in its place.

\begin{example}{Lie Algebras}{Lie Algebras}
  \begin{itemize}
    \item The vector space $\mathbb{R}^3$ with the cross product as the Lie bracket is a Lie algebra. The cross product is bilinear, antisymmetric, and satisfies the Jacobi identity.
    \item Let $A$ be any associative algebra over $\mathbb{R}$ or $\mathbb{C}$. Then $A$ can be made into a Lie algebra by defining the Lie bracket as
      \begin{equation}
        [x,y]=xy-yx
      \end{equation}
      for all $x,y\in A$. Recall that an algebra is a vector space equipped with a bilinear multiplication operation.

      Expanding the Jacobi identity gives pairs of terms such as $x(yz)$ and $(xy)z$, which cancel by associativity.
  \end{itemize}
\end{example}

\begin{remark}
  We shall prove later that every Lie algebra can be embedded into an associative algebra in such a way that the Lie bracket is given by the commutator of the associative multiplication.
\end{remark}

\begin{example}{$sl(n;\mathbb{C})$}{slnmathbbC}
  Let $sl(n;\mathbb{C})$ be the set of all $X \in M_n(\mathbb{C})$ such that $\tr(X)=0$. Then $sl(n;\mathbb{C})$ is a Lie algebra with the Lie bracket given by the commutator $[X,Y]=XY-YX$.

  \begin{proof}
    For any $X,Y\in M_n(\mathbb{C})$, we have
    \begin{equation*}
      \tr([X,Y])=\tr(XY-YX)=\tr(XY)-\tr(YX)=0,
    \end{equation*}
    so $[\cdot,\cdot]$ is closed on $sl(n;\mathbb{C})$.
  \end{proof}
\end{example}

\begin{definition}{Subalgebra of Lie Algebras}{Subalgebra of Lie Algebras}
  A \textbf{subalgebra} of a Lie algebra $\mathfrak{g}$ is a subspace $\mathfrak{h}\subseteq\mathfrak{g}$ such that $[\cdot,\cdot]$ is closed on $\mathfrak{h}$.

  A subalgebra $\mathfrak{h}$ of a Lie algebra $\mathfrak{g}$ is called an \textbf{ideal} if $[x,h]\in\mathfrak{h}$ for all $x\in\mathfrak{g}$ and $h\in\mathfrak{h}$.

  The \textbf{center} of a Lie algebra $\mathfrak{g}$ is the set of all elements $z\in\mathfrak{g}$ such that $[z,x]=0$ for all $x\in\mathfrak{g}$.
\end{definition}


The distinction is which inputs are allowed: a subalgebra requires $[\mathfrak h,\mathfrak h]\subseteq\mathfrak h$, while an ideal requires $[\mathfrak g,\mathfrak h]\subseteq\mathfrak h$. Figure~\ref{fig:subalgebra-versus-ideal} gives a concrete example.

\begin{figure}[H]
    \centering
    \incfig{subalgebra-versus-ideal}
    \caption{In $(\mathbb R^3,\times)$, the line $\mathbb Re_3$ is a subalgebra because brackets of vectors on the line vanish. It is not an ideal: taking a bracket with the outside direction $e_1$ gives $-e_2$, outside the line.}
    \label{fig:subalgebra-versus-ideal}
\end{figure}

\begin{definition}{Homomorphisms}{Homomorphisms}
  If $\mathfrak{g}$ and $\mathfrak{h}$ are Lie algebras, a \textbf{homomorphism} from $\mathfrak{g}$ to $\mathfrak{h}$ is a linear map $\phi:\mathfrak{g}\to\mathfrak{h}$ such that
  \begin{equation}
    \phi([x,y])=[\phi(x),\phi(y)]
  \end{equation}
  for all $x,y\in\mathfrak{g}$. If $\phi$ is bijective, it is called an \textbf{isomorphism}. If $\mathfrak{g}=\mathfrak{h}$, then $\phi$ is called an \textbf{endomorphism}, and if it is bijective, it is called an \textbf{automorphism}.
\end{definition}

\begin{definition}{Adjoint Map}{Adjoint Map}
  If $\mathfrak{g}$ is a Lie algebra, and $x \in \mathfrak{g}$, we define the \textbf{adjoint map} $\ad_x:\mathfrak{g}\to\mathfrak{g}$ by
  \begin{equation}
    \ad_x(y)=[x,y]
  \end{equation}
  which is also called the \textbf{adjoint representation} of $\mathfrak{g}$.
\end{definition}
The Jacobi identity can be rewritten in terms of the adjoint map as
\begin{equation}
  \ad_x([y,z])=[\ad_x(y),z]+[y,\ad_x(z)].
\end{equation}

The assignment $x\mapsto\ad_x$ is a Lie algebra homomorphism into $\mathfrak{gl}(\mathfrak{g})$, as the following identity shows.

\begin{proposition}{}{}
  If $\mathfrak{g}$ is a Lie algebra, then we have
  \begin{equation}
    \ad_{[x,y]}=[\ad_x,\ad_y]
  \end{equation}
\end{proposition}
\begin{proof}
  Applying both sides to $z\in\mathfrak{g}$ gives the Jacobi identity.
\end{proof}

\begin{definition}{Direct Sums of Lie Algebras}{Direct Sums of Lie Algebras}
  If $\mathfrak{g}_1$ and $\mathfrak{g}_2$ are Lie algebras, then the \textbf{direct sum} $\mathfrak{g}_1\oplus\mathfrak{g}_2$ is the vector space $\mathfrak{g}_1\oplus\mathfrak{g}_2$ with the Lie bracket defined by
  \begin{equation}
    [(x_1,x_2),(y_1,y_2)]=([x_1,y_1],[x_2,y_2])
  \end{equation}
  for all $(x_1,x_2),(y_1,y_2)\in\mathfrak{g}_1\oplus\mathfrak{g}_2$.
\end{definition}

We can describe a Lie algebra in terms of a basis:
\begin{definition}{Structure Constants}{Structure Constants}
  Let $\mathfrak{g}$ be a Lie algebra with basis $\{e_1,\ldots,e_n\}$. Then we can write
  \begin{equation}
    [e_i,e_j]=\sum_{k=1}^n c_{ij}^k e_k
  \end{equation}
  for some constants $c_{ij}^k\in\mathbb{R}$ or $\mathbb{C}$, called the \textbf{structure constants} of $\mathfrak{g}$ with respect to the basis $\{e_1,\ldots,e_n\}$.
\end{definition}

\section{Simple, Solvable, and Nilpotent Lie Algebras}
In this section, we consider various special types of Lie algebras.

\begin{definition}{Simple}{Simple}
  A Lie algebra $\mathfrak{g}$ is called \textbf{irreducible} if it has no nontrivial ideals, i.e. the only ideals of $\mathfrak{g}$ are $\{0\}$ and $\mathfrak{g}$ itself. A Lie algebra $\mathfrak{g}$ is called \textbf{simple} if it is irreducible and $\dim\mathfrak{g}>2$.
\end{definition}

A one-dimensional Lie algebra is irreducible, since it has no nontrivial subspaces and hence no nontrivial ideals. Nevertheless, such a Lie algebra is, by definition, not considered simple.

A one-dimensional Lie algebra $\mathfrak{g}$ is necessarily commutative, since $[ax, bx]=ab[x,x]=0$ for all $a,b\in\mathbb{R}$ or $\mathbb{C}$ and $x\in\mathfrak{g}$. On the other hand, if $\mathfrak{g}$ is commutative, then every subspace of $\mathfrak{g}$ is an ideal, so in this case $\mathfrak{g}$ is irreducible if and only if it is one-dimensional. Thus, an equivalent definition of simple is that a Lie algebra is simple if it is \textbf{noncommutative and irreducible}.

\begin{remark}
  There is an analogy between groups and Lie algebras, in which the role of subgroups is played by subalgebras and the role of normal subgroups is played by ideals. For example, the kernel of a Lie algebra homomorphism is always an ideal, just as the kernel of a Lie group homomorphism is always a normal subgroup.

  In group theory, the word ``simple'' is used to describe a group with no nontrivial normal subgroups, since normal subgroups are the ones that can be used to form quotient groups, and thus ``factor'' the group.
\end{remark}

\begin{proposition}{}{}
  The Lie algebra $sl(2;\mathbb{C})$ is simple.
\end{proposition}
\begin{proof}
  Consider the basis
  \begin{equation*}
    X = \begin{pmatrix}
      0 & 1 \\
      0 & 0
    \end{pmatrix},\quad
    Y = \begin{pmatrix}
      0 & 0 \\
      1 & 0
    \end{pmatrix},\quad
    H = \begin{pmatrix}
      1 & 0 \\
      0 & -1
    \end{pmatrix}
  \end{equation*}
  Let $\mathfrak{h}$ be a nonzero ideal and choose $0\neq Z=aX+bH+cY\in\mathfrak{h}$. We show that $\mathfrak{h}$ contains all three basis elements. First, we compute the commutators:
  \begin{equation*}
    [X, Y] = H,\quad [H, X] = 2X,\quad [H, Y] = -2Y
  \end{equation*}
  and we have
  \begin{equation*}
    [X, Z] = -2bX + cH, \quad [Y, Z] = -aH + 2bY, \quad [H, Z] = 2aX - 2cY
  \end{equation*}
  If $c \neq 0$, then $[X, [X, Z]] = -2cX \in \mathfrak{h}$, so $X \in \mathfrak{h}$. Then $[X, Z] = -2bX + cH \in \mathfrak{h}$, so $H \in \mathfrak{h}$, and $Y \in \mathfrak{h}$ as well. If $a\neq0$, the analogous argument starts with $[Y,[Y,Z]]=-2aY$. If $a=c=0$, then $b\neq0$ and $Z=bH$ already yields $X$ and $Y$ by taking brackets. Thus every nonzero ideal is the whole algebra, and $sl(2;\mathbb{C})$ is simple.
\end{proof}

\begin{definition}{Commutator Ideal}{Commutator Ideal}
  If $\mathfrak{g}$ is a Lie algebra, then the \textbf{commutator ideal} of $\mathfrak{g}$ is denoted by $[\mathfrak{g},\mathfrak{g}]$ and is defined as the space of linear combinations of commutators, which is spanned by all elements of the form $[x,y]$ for $x,y\in\mathfrak{g}$.
\end{definition}

\begin{definition}{Solvable}{Solvable}
  For any Lie algebra $\mathfrak{g}$, we define the \textbf{derived series} of $\mathfrak{g}$ as follows:
  \begin{equation}
    \mathfrak{g}_0 = \mathfrak{g},\quad \mathfrak{g}_{k+1} = [\mathfrak{g}_k,\mathfrak{g}_k]
  \end{equation}
  These subalgebras are called the \textbf{derived subalgebras} of $\mathfrak{g}$. A Lie algebra $\mathfrak{g}$ is called \textbf{solvable} if there exists an integer $n$ such that $\mathfrak{g}_n = \{0\}$.
\end{definition}

The lower central series descends more slowly:
\begin{definition}{Nilpotent}{Nilpotent}
  For any Lie algebra $\mathfrak{g}$, we define the \textbf{lower central series} of $\mathfrak{g}$ as follows: Let $\mathfrak{g}^0 = \mathfrak{g}$, and $\mathfrak{g}^{k+1} = [\mathfrak{g},\mathfrak{g}^k]$ for $k \geq 0$. A Lie algebra $\mathfrak{g}$ is called \textbf{nilpotent} if there exists an integer $n$ such that $\mathfrak{g}^n = \{0\}$.
\end{definition}

It is not difficult to show that $\mathfrak{g}_k \subseteq \mathfrak{g}^k$ for all $k$, so every nilpotent Lie algebra is solvable.

\begin{remark}
  By bilinearity, commutators of basis elements suffice to compute the derived and lower central series. When each such commutator is zero or a scalar multiple of a basis element, these computations reduce to tracking the basis elements that remain at each stage.
\end{remark}

\begin{example}{Upper Triangular Matrices}{Upper Triangular Matrices}
  Let $\mathfrak{g}$ be the Lie algebra of all $n \times n$ strictly upper triangular matrices (with zero diagonal). Then $\mathfrak{g}$ is easily seen to be nilpotent, since successive products move nonzero entries farther above the diagonal.

  \medskip
  If $\mathfrak{g} \subseteq M_2(\mathbb{C})$ is the Lie algebra of all $2 \times 2$ upper triangular matrices (with possibly nonzero diagonal), then $\mathfrak{g}$ is solvable but not nilpotent.
\end{example}


Figure~\ref{fig:derived-and-lower-central-series} contrasts the two series. In the first example, $[P,Q]=Z$ and $Z$ is central. In the second, take $D=\operatorname{diag}(1,0)$ and $N=E_{12}$; then $[D,N]=N$, so bracketing with the original algebra keeps producing $N$. Angle brackets denote linear spans.

\begin{figure}[H]
    \centering
    \incfig{derived-and-lower-central-series}
    \caption{The derived series brackets the current space with itself; the lower central series brackets it with the original algebra. Both reach zero for the Heisenberg algebra, but only the derived series reaches zero for upper triangular $2\times2$ matrices.}
    \label{fig:derived-and-lower-central-series}
\end{figure}

\section{The Lie Algebra of a Matrix Lie Group}
In this section, we associate to each matrix Lie group $G$ a Lie algebra $\mathfrak{g}$.

\begin{definition}{Lie Algebra of a Matrix Lie Group}{Lie Algebra of a Matrix Lie Group}
  Let $G$ be a matrix Lie group. The \textbf{Lie algebra} of $G$ is the set
  \begin{equation}
    \mathfrak{g} = \{X \in M_n(\mathbb{C}) : e^{tX} \in G \text{ for all } t \in \mathbb{R}\}.
  \end{equation}
\end{definition}
Thus $\mathfrak{g}$ consists of the generators of one-parameter subgroups of $G$.

\begin{remark}
  We require $e^{tX} \in G$ for all $t \in \mathbb{R}$, not for all $t \in \mathbb{C}$.
\end{remark}

Since $t\mapsto e^{tX}$ is a continuous path through the identity, it stays in the identity component of $G$. Hence $G$ and its identity component have the same Lie algebra.

\begin{theorem}{}{}
  Let $G$ be a matrix Lie group with Lie algebra $\mathfrak{g}$. If $X,Y \in \mathfrak{g}$, then
  \begin{itemize}
    \item $AXA^{-1} \in \mathfrak{g}$ for all $A \in G$,
    \item $sX \in \mathfrak{g}$ for all $s \in \mathbb{R}$,
    \item $X + Y \in \mathfrak{g}$,
    \item $[X,Y] \in \mathfrak{g}$.
  \end{itemize}
\end{theorem}
\begin{proof}
  Conjugation preserves the group and therefore its infinitesimal generators. Indeed,
  \begin{equation*}
    e^{t(AXA^{-1})}=Ae^{tX}A^{-1}\in G
    \qquad\text{for all }t\in\mathbb{R},
  \end{equation*}
  so $AXA^{-1}\in\mathfrak{g}$.

  For scalar multiplication, $e^{t(sX)} = e^{(ts)X} \in G$ for all $t \in \mathbb{R}$, so $sX \in \mathfrak{g}$.

  For addition, the Lie product formula gives $e^{t(X+Y)} = \lim_{n \to \infty} (e^{tX/n} e^{tY/n})^n \in G$ for all $t \in \mathbb{R}$. Each term belongs to $G$, and the limit is invertible. Since $G$ is closed in $GL(n;\mathbb{C})$, the limit also belongs to $G$, so $X+Y\in\mathfrak{g}$.

  For closure under the bracket, we have
  \begin{equation}
    \frac{\mathrm{d}}{\mathrm{d} t} (e^{tX} Y e^{-tX}) \bigg|_{t=0} = XY - YX = [X,Y] = \lim_{h \to 0} \frac{e^{hX} Y e^{-hX} - Y}{h}
  \end{equation}
  and $e^{tX} Y e^{-tX} \in \mathfrak{g}$ for all $t \in \mathbb{R}$, so $[X,Y] \in \mathfrak{g}$ because $\mathfrak{g}$ is a finite-dimensional real vector subspace and hence is closed in $M_n(\mathbb{C})$.
\end{proof}

\begin{remark}
  Note that even if the elements of $G$ have complex entries, the Lie algebra $\mathfrak{g}$ is not necessarily a complex vector space, since the scalar-multiplication statement only guarantees that $sX \in \mathfrak{g}$ for all $s \in \mathbb{R}$, not for all $s \in \mathbb{C}$.
\end{remark}

\begin{definition}{Complex Lie Group}{Complex Lie Group}
  A matrix Lie group $G$ is called a \textbf{complex Lie group} if its Lie algebra $\mathfrak{g}$ is a complex subspace of $M_n(\mathbb{C})$. That is, if $i X \in \mathfrak{g}$ for all $X \in \mathfrak{g}$.
\end{definition}

\begin{proposition}{}{}
  If $G$ is commutative, then its Lie algebra $\mathfrak{g}$ is also commutative.
\end{proposition}
\begin{proof}
  For any $X, Y \in M_n(\mathbb{C})$, we have
  \begin{equation}
    [X, Y] = \frac{\mathrm{d}}{\mathrm{d} t} \left( \frac{\mathrm{d}}{\mathrm{d} s} (e^{tX} e^{sY} e^{-tX}) \bigg|_{s=0} \right) \bigg|_{t=0}
  \end{equation}
  If $G$ is commutative, then $e^{tX} e^{sY} e^{-tX} = e^{sY}$ for all $t, s \in \mathbb{R}$, so the right-hand side is zero, and thus $[X, Y] = 0$ for all $X, Y \in \mathfrak{g}$.
\end{proof}

\section{Examples}
\begin{proposition}{$GL$ and $SL$}{GL and SL}
  The Lie algebra of $GL(n;\mathbb{C})$ is $M_n(\mathbb{C})$.

  The Lie algebra of $GL(n;\mathbb{R})$ is $M_n(\mathbb{R})$.

  The Lie algebra of $SL(n;\mathbb{C})$ consists of all $X \in M_n(\mathbb{C})$ such that $\tr(X) = 0$, i.e. $sl(n;\mathbb{C})$.

  The Lie algebra of $SL(n;\mathbb{R})$ consists of all $X \in M_n(\mathbb{R})$ such that $\tr(X) = 0$, i.e. $sl(n;\mathbb{R})$.
\end{proposition}
\begin{proof}
  If $X \in M_n(\mathbb{C})$, then $e^{tX} \in GL(n;\mathbb{C})$ for all $t \in \mathbb{R}$, so $M_n(\mathbb{C}) \subseteq \mathfrak{g} \subseteq M_n(\mathbb{C})$, and thus $\mathfrak{g} = M_n(\mathbb{C})$. The same argument works for $GL(n;\mathbb{R})$. The $SL(n;\mathbb{C})$ and $SL(n;\mathbb{R})$ cases follow from the fact that $\det(e^{tX}) = e^{t \tr(X)}$.
\end{proof}

\begin{proposition}{$U$ and $O$}{U and O}
  The Lie algebra of $U(n)$ consists of all $X \in M_n(\mathbb{C})$ such that $X + X^* = 0$, i.e. the set of all skew-Hermitian matrices.

  The Lie algebra of $SU(n)$ consists of all $X \in M_n(\mathbb{C})$ such that $X + X^* = 0$ and $\tr(X) = 0$, i.e. the set of all skew-Hermitian matrices with trace zero.

  The Lie algebra of $O(n)$ consists of all $X \in M_n(\mathbb{R})$ such that $X + X^T = 0$, i.e. the set of all skew-symmetric matrices.

  The Lie algebra of $SO(n)$ is the same as that of $O(n)$, since the determinant condition is automatically satisfied for skew-symmetric matrices.

  We denote the corresponding Lie algebras by $\mathfrak{u}(n)$, $\mathfrak{su}(n)$, and $\mathfrak{so}(n)$, respectively.
\end{proposition}
\begin{proof}
  The identity $(e^{tX})^*=e^{tX^*}=e^{-tX}$ holds for all real $t$ exactly when $X+X^*=0$, as follows by differentiating at $t=0$ and exponentiating the resulting condition. Requiring $\det(e^{tX})=1$ for all $t$ additionally gives $\tr(X)=0$. The orthogonal cases follow similarly, using the transpose in place of the adjoint.
\end{proof}

\begin{proposition}{$O(n;k)$, $SO(n;k)$, and $Sp$}{Onk, SOnk, and Sp}
  The Lie algebra of $O(n;k)$ consists of all $X \in M_{n+k}(\mathbb{R})$ such that $X^T I_{n,k} + I_{n,k} X = 0$, i.e. the set of all matrices $X$ such that $I_{n,k} X$ is skew-symmetric.

  The Lie algebra of $SO(n;k)$ is the same as that of $O(n;k)$, since the determinant condition is automatically satisfied for matrices satisfying the above condition.

  The Lie algebra of $Sp(n;\mathbb{R})$ consists of all $X \in M_{2n}(\mathbb{R})$ such that $X^T J + J X = 0$, i.e. the set of all matrices $X$ such that $J X$ is symmetric.

  The Lie algebra of $Sp(n;\mathbb{C})$ consists of all $X \in M_{2n}(\mathbb{C})$ such that $X^T J + J X = 0$, i.e. the set of all matrices $X$ such that $J X$ is symmetric.

  We denote the corresponding Lie algebras by $\mathfrak{so}(n;k)$, and $\mathfrak{sp}(n;\mathbb{R})$, and $\mathfrak{sp}(n;\mathbb{C})$, respectively.
\end{proposition}

\begin{proposition}{Heisenberg Group}{Heisenberg Group}
  The Heisenberg group $H_n$ is the set of all $(n+2) \times (n+2)$ matrices of the form
  \begin{equation}
    \begin{pmatrix}
      1 & x & z \\
      0 & I_n & y \\
      0 & 0 & 1
    \end{pmatrix}
  \end{equation}
  and its Lie algebra $\mathfrak{h}_n$ is the set of all $(n+2) \times (n+2)$ matrices of the form
  \begin{equation}
    \begin{pmatrix}
      0 & x & z \\
      0 & 0 & y \\
      0 & 0 & 0
    \end{pmatrix}
  \end{equation}
\end{proposition}

\begin{example}{Pauli Matrices}{Pauli Matrices}
  The following elements form a basis for the Lie algebra $\mathfrak{su}(2)$:
  \begin{equation}
    E_1 = \frac{1}{2} \begin{pmatrix}
      i & 0 \\
      0 & -i
    \end{pmatrix},\quad
    E_2 = \frac{1}{2} \begin{pmatrix}
      0 & i \\
      i & 0
    \end{pmatrix},\quad
    E_3 = \frac{1}{2} \begin{pmatrix}
      0 & -1 \\
      1 & 0
    \end{pmatrix}
  \end{equation}
  These satisfy $[E_1, E_2] = E_3$, $[E_2, E_3] = E_1$, and $[E_3, E_1] = E_2$.

  The following elements form a basis for the Lie algebra $\mathfrak{so}(3)$:
  \begin{equation}
    F_1 = \begin{pmatrix}
      0 & 0 & 0 \\
      0 & 0 & -1 \\
      0 & 1 & 0
    \end{pmatrix},\quad
    F_2 = \begin{pmatrix}
      0 & 0 & 1 \\
      0 & 0 & 0 \\
      -1 & 0 & 0
    \end{pmatrix},\quad
    F_3 = \begin{pmatrix}
      0 & -1 & 0 \\
      1 & 0 & 0 \\
      0 & 0 & 0
    \end{pmatrix}
  \end{equation}
  These satisfy $[F_1, F_2] = F_3$, $[F_2, F_3] = F_1$, and $[F_3, F_1] = F_2$.

  Since the two bases have the same commutation relations, $\mathfrak{su}(2)$ and $\mathfrak{so}(3)$ are isomorphic as Lie algebras.
\end{example}

\section{Lie Group and Lie Algebra Homomorphisms}
We now establish the correspondence between group homomorphisms and their infinitesimal maps, introduced in the geometric overview.

\begin{theorem}{Lie Group and Lie Algebra Homomorphisms}{Lie Group and Lie Algebra Homomorphisms}
  Let $G$ and $H$ be matrix Lie groups with Lie algebras $\mathfrak{g}$ and $\mathfrak{h}$, respectively. If $\Phi:G\to H$ is a Lie group homomorphism, then there exists a unique real-linear map $\phi:\mathfrak{g}\to\mathfrak{h}$ such that
  \begin{equation}
    \Phi(e^{X}) = e^{\phi(X)}
  \end{equation}
  for all $X \in \mathfrak{g}$. The map $\phi$ has the following properties:
  \begin{itemize}
    \item $\phi(A X A^{-1}) = \Phi(A) \phi(X) \Phi(A)^{-1}$ for all $A \in G$ and $X \in \mathfrak{g}$,
    \item $\phi([X,Y]) = [\phi(X),\phi(Y)]$ for all $X,Y \in \mathfrak{g}$. This means $\phi$ is a Lie algebra homomorphism.
    \item We have $\displaystyle \phi(X) = \frac{\mathrm{d}}{\mathrm{d} t} \Phi(e^{tX}) \bigg|_{t=0}$ for all $X \in \mathfrak{g}$.
  \end{itemize}
\end{theorem}


Figure~\ref{fig:homomorphisms-and-exponentials} organizes the preceding theorem into a commuting diagram. The upper row lies in Lie algebras and the lower row in Lie groups.

\begin{figure}[H]
    \centering
    \incfig{homomorphisms-and-exponentials}
    \caption{Exponentiate and then apply $\Phi$, or apply its differential $\phi$ and then exponentiate: both routes give $\Phi(e^X)=e^{\phi(X)}$. The group map acts on finite transformations; its differential acts on their generators.}
    \label{fig:homomorphisms-and-exponentials}
\end{figure}

\begin{example}{}{}
  Let $\Phi: SU(2) \to SO(3)$ be the Lie group homomorphism defined by \ref{exp:SU2 and SO3}. Then the associated Lie algebra homomorphism $\phi: \mathfrak{su}(2) \to \mathfrak{so}(3)$ is given by $\phi(E_i) = F_i$, which is a Lie algebra isomorphism.
\end{example}

\begin{proposition}{}{}
  If $\Phi: G \to H$ is a Lie group homomorphism and $\phi: \mathfrak{g} \to \mathfrak{h}$ is the associated Lie algebra homomorphism, then the kernel of $\Phi$ is a closed, normal subgroup of $G$, and the Lie algebra of $\ker(\Phi)$ is the kernel of $\phi$.
\end{proposition}

We now return to the group adjoint action from the geometric overview. Its differential will recover the algebraic map $\ad$ defined in \ref{def:Adjoint Map}.

\begin{definition}{Adjoint Representation of a Lie Group}{Adjoint Representation of a Lie Group}
  Let $G$ be a matrix Lie group with Lie algebra $\mathfrak{g}$. For each $A \in G$, we define the \textbf{adjoint representation} $\Ad_A: \mathfrak{g} \to \mathfrak{g}$ by
  \begin{equation}
    \Ad_A(X) = A X A^{-1}
  \end{equation}
  for all $X \in \mathfrak{g}$.
\end{definition}

\begin{proposition}{}{}
  Let $G$ be a matrix Lie group with Lie algebra $\mathfrak{g}$. Let $GL(\mathfrak{g})$ be the group of all invertible linear transformations on $\mathfrak{g}$. Then the map $\Ad: G \to GL(\mathfrak{g})$ defined by $A \mapsto \Ad_A$ is a Lie group homomorphism.
\end{proposition}

\begin{proof}
  First, $\Ad_A$ is a linear transformation on $\mathfrak{g}$, since for any $X,Y\in\mathfrak{g}$ and $\alpha,\beta\in\mathbb{R}$, we have
  \begin{equation*}
    \Ad_A(\alpha X+\beta Y)=\alpha\Ad_A(X)+\beta\Ad_A(Y).
  \end{equation*}
  Moreover, $\Ad_A$ is invertible, with inverse $\Ad_{A^{-1}}$. Thus, $\Ad_A\in GL(\mathfrak{g})$.

  For any $A,B\in G$ and $X\in\mathfrak{g}$, we have
  \begin{equation*}
    \begin{aligned}
      \Ad_{AB}(X)&=(AB)X(AB)^{-1}\\
      &=A(BXB^{-1})A^{-1}\\
      &=\Ad_A(\Ad_B(X)).
    \end{aligned}
  \end{equation*}
  Hence $\Ad_{AB}=\Ad_A\circ\Ad_B$, so $\Ad$ is a group homomorphism. The continuity of $\Ad$ follows from the continuity of matrix multiplication and inversion.
\end{proof}

Since $\Ad:G\to GL(\mathfrak{g})$ is a Lie group homomorphism, it induces a Lie algebra homomorphism from $\mathfrak{g}$ to the Lie algebra of $GL(\mathfrak{g})$, denoted by $\mathfrak{gl}(\mathfrak{g})$. Here, $\mathfrak{gl}(\mathfrak{g})$ is the space of all linear transformations on $\mathfrak{g}$, equipped with the commutator bracket. We denote this associated Lie algebra homomorphism by
\begin{equation}
  \ad:\mathfrak{g}\to\mathfrak{gl}(\mathfrak{g}).
\end{equation}
By the preceding theorem, we have
\begin{equation}
  e^{\ad_X}=\Ad_{e^X}
\end{equation}
for all $X\in\mathfrak{g}$.

\begin{proposition}{}{}
  Let $G$ be a matrix Lie group with Lie algebra $\mathfrak{g}$, and let $\ad:\mathfrak{g}\to\mathfrak{gl}(\mathfrak{g})$ be the Lie algebra homomorphism associated with $\Ad:G\to GL(\mathfrak{g})$. Then for all $X,Y\in\mathfrak{g}$, we have
  \begin{equation}
    \ad_X(Y)=[X,Y].
  \end{equation}
\end{proposition}
\begin{proof}
  By the preceding theorem, we have
  \begin{equation*}
    \ad_X=\frac{\mathrm{d}}{\mathrm{d}t}\Ad_{e^{tX}}\bigg|_{t=0}.
  \end{equation*}
  Thus, for any $Y\in\mathfrak{g}$,
  \begin{equation*}
    \begin{aligned}
      \ad_X(Y)
      &=\frac{\mathrm{d}}{\mathrm{d}t}\left(e^{tX}Ye^{-tX}\right)\bigg|_{t=0}\\
      &=XY-YX\\
      &=[X,Y].
    \end{aligned}
  \end{equation*}
\end{proof}

\begin{remark}
  This proposition shows that the adjoint representation of a Lie algebra defined in \ref{def:Adjoint Map} is precisely the Lie algebra homomorphism associated with the adjoint representation of its Lie group.

  In particular, we should distinguish the two maps
  \begin{equation}
    \Ad:G\to GL(\mathfrak{g}),\quad
    \ad:\mathfrak{g}\to\mathfrak{gl}(\mathfrak{g}).
  \end{equation}
  The former describes the action of conjugation on the Lie algebra, while the latter describes its infinitesimal action, given by the Lie bracket.
\end{remark}

As a consequence, we obtain the following useful identity, which relates conjugation by the matrix exponential to the exponential of the adjoint map.

\begin{proposition}{}{}
  For any $X,Y\in M_n(\mathbb{C})$, let $\ad_X:M_n(\mathbb{C})\to M_n(\mathbb{C})$ be given by $\ad_X(Y)=[X,Y]$. Then
  \begin{equation}
    e^XYe^{-X}=\Ad_{e^X}(Y)=e^{\ad_X}(Y),
  \end{equation}
  where
  \begin{equation}
    e^{\ad_X}(Y)
    =Y+[X,Y]+\frac{1}{2!}[X,[X,Y]]
    +\frac{1}{3!}[X,[X,[X,Y]]]+\cdots.
  \end{equation}
\end{proposition}
\begin{proof}
  Since $GL(n;\mathbb{C})$ is a matrix Lie group with Lie algebra $M_n(\mathbb{C})$, the preceding results give
  \begin{equation*}
    \Ad_{e^X}=e^{\ad_X}.
  \end{equation*}
  Applying both sides to $Y$, we obtain
  \begin{equation*}
    e^XYe^{-X}=e^{\ad_X}(Y).
  \end{equation*}
  Expanding the exponential of the linear transformation $\ad_X$, we have
  \begin{equation*}
    \begin{aligned}
      e^{\ad_X}(Y)
      &=\sum_{k=0}^{\infty}\frac{1}{k!}\ad_X^k(Y)\\
      &=Y+[X,Y]+\frac{1}{2!}[X,[X,Y]]+\cdots.
    \end{aligned}
  \end{equation*}
\end{proof}

\begin{remark}
  The identity $e^XYe^{-X}=e^{\ad_X}(Y)$ shows that conjugation by $e^X$ can be expressed entirely in terms of the Lie bracket. Thus, the adjoint representation provides a direct connection between the group operation and the algebraic structure of its Lie algebra.
\end{remark}

\begin{remark}
  To distinguish the tangent spaces involved, consider conjugation as a two-parameter family. Let $X,Y\in\mathfrak{g}$, and consider the two one-parameter subgroups $A(t)=e^{tX}$ and $B(s)=e^{sY}$. Conjugation gives
  \begin{equation}
    F(t,s)=e^{tX}e^{sY}e^{-tX}\in G.
  \end{equation}
  We can differentiate this expression with respect to either $s$ or $t$, which have different geometric meanings.

  First, fixing $t$ and differentiating with respect to $s$ at $s=0$, we obtain
  \begin{equation}
    \frac{\partial F}{\partial s}\bigg|_{s=0}
    =e^{tX}Ye^{-tX}
    =\Ad_{e^{tX}}(Y).
  \end{equation}
  Since $F(t,0)=I$, this derivative belongs to $T_I G=\mathfrak{g}$. Thus, $\Ad$ describes how conjugation transforms tangent vectors at the identity.

  On the other hand, fixing $s$ and differentiating with respect to $t$ at $t=0$, we obtain
  \begin{equation}
    \frac{\partial F}{\partial t}\bigg|_{t=0}
    =Xe^{sY}-e^{sY}X.
  \end{equation}
  Since $F(0,s)=e^{sY}$, this derivative belongs to $T_{e^{sY}}G$, rather than necessarily to $\mathfrak{g}$. It describes the infinitesimal change of the group element $e^{sY}$ under conjugation generated by $X$.

  Finally, differentiating with respect to both parameters at the identity, we have
  \begin{equation}
    \begin{aligned}
      \frac{\partial^2 F}{\partial t\partial s}\bigg|_{t=s=0}
      &=\frac{\mathrm{d}}{\mathrm{d}t}
      \left(e^{tX}Ye^{-tX}\right)\bigg|_{t=0}\\
      &=XY-YX\\
      &=[X,Y].
    \end{aligned}
  \end{equation}
  The same result is obtained if we differentiate first with respect to $t$ and then with respect to $s$, since the mixed partial derivatives commute.

  The derivatives and their tangent spaces are:
  \begin{center}
    \renewcommand{\arraystretch}{1.5}
    \begin{tabular}{|c|c|c|}
      \hline
      \textbf{Operation} & \textbf{Result} & \textbf{Tangent Space}\\
      \hline
      $\displaystyle\frac{\partial F}{\partial s}\bigg|_{s=0}$
      & $\Ad_{e^{tX}}(Y)$
      & $\mathfrak{g}=T_I G$\\
      \hline
      $\displaystyle\frac{\partial F}{\partial t}\bigg|_{t=0}$
      & $Xe^{sY}-e^{sY}X$
      & $T_{e^{sY}}G$\\
      \hline
      $\displaystyle\frac{\partial^2 F}{\partial t\partial s}\bigg|_{t=s=0}$
      & $\ad_X(Y)=[X,Y]$
      & $\mathfrak{g}=T_I G$\\
      \hline
    \end{tabular}
  \end{center}

  Therefore, the Lie bracket can be interpreted as the mixed infinitesimal variation of conjugation in two directions. More precisely, the group commutator satisfies
  \begin{equation}
    e^{tX}e^{sY}e^{-tX}e^{-sY}
    =I+ts[X,Y]+O(t^2s+ts^2).
  \end{equation}
  Hence $[X,Y]$ measures the leading-order failure of the two one-parameter subgroups to commute. In this sense, $\Ad$ describes finite conjugation on the Lie algebra, while $\ad$ describes its infinitesimal behavior.
\end{remark}

Figure~\ref{fig:adjoint-conjugation-geometry} illustrates these derivatives using
\begin{equation*}
  X=\begin{pmatrix}0&-1\\1&0\end{pmatrix},\qquad
  Y=\begin{pmatrix}0&1\\1&0\end{pmatrix}
  \quad\text{in }\mathfrak{sl}(2,\mathbb{R}).
\end{equation*}
Here $F(t,s)=\cosh s\,I+\sinh s\,\Ad_{e^{tX}}Y$ lies on the surface $x^2-z^2-w^2=1$, $x>0$, in the slice $A=xI+zY+w\operatorname{diag}(1,-1)$. The blue curves share the base point $I$, whereas the red conjugation orbit passes through $e^{s_0Y}$. Differentiating the blue curves at $I$ gives the vectors in the right panel; the red arrow there shows their variation at $t=0$, based at $Y$ in the vector space $\mathfrak{g}$.

\begin{figure}[H]
    \centering
    \incfig{adjoint-conjugation-geometry}
    \caption{Conjugation and its two infinitesimal variations. Left: blue tangent arrows at $I$ represent $\partial_sF(t,0)$; the red tangent arrow at $e^{s_0Y}$ represents $\partial_tF(0,s_0)$. Right: the blue vectors are $\Ad_{e^{tX}}Y$, and their derivative at $t=0$ is $[X,Y]$. The panels show two-dimensional slices.}
    \label{fig:adjoint-conjugation-geometry}
\end{figure}

\section{The Complexification of a Real Lie Algebra}
The Lie algebra of a matrix Lie group is always a real vector space. To study it over $\mathbb{C}$, we pass to its \textbf{complexification}.

\begin{definition}{Complexification}{Complexification}
  If $V$ is a finite-dimensional real vector space, then the \textbf{complexification} of $V$ is the space of formal linear combinations $v_1 + i v_2$ with $v_1, v_2 \in V$, denoted by $V_\mathbb{C}$.
\end{definition}


For the simplest example, $\mathbb R_{\mathbb C}\cong\mathbb C$, as in Figure~\ref{fig:complexification-real-and-imaginary}. More generally, a real basis of $V$ becomes a complex basis of $V_{\mathbb C}$. The distinction between this formal construction and an ambient complex span is explained below.

\begin{figure}[H]
    \centering
    \incfig{complexification-real-and-imaginary}
    \caption{Complexifying the real line adds a formal imaginary direction: $a+ib$ records two real components but is one complex coordinate. This is an illustration for $V=\mathbb R$; the same construction applies to each vector-space direction.}
    \label{fig:complexification-real-and-imaginary}
\end{figure}

We shall see that there is a natural way to complexify a real Lie algebra $\mathfrak{g}$ to obtain a complex Lie algebra $\mathfrak{g}_\mathbb{C}$.
\begin{proposition}{Complexification of Lie Algebra}{Complexification of Lie Algebra}
  Let $\mathfrak{g}$ be a finite-dimensional real Lie algebra, and $\mathfrak{g}_\mathbb{C}$ its complexification. Then the bracket operation on $\mathfrak{g}$ has a unique extension to a bracket operation on $\mathfrak{g}_\mathbb{C}$, making $\mathfrak{g}_\mathbb{C}$ into a complex Lie algebra.
\end{proposition}
\begin{proof}
  Complex bilinearity forces the formula
  \begin{equation*}
    [X_1 + i Y_1, X_2 + i Y_2] = [X_1, X_2] - [Y_1, Y_2] + i([X_1, Y_2] + [Y_1, X_2])
  \end{equation*}
  This formula defines the unique complex-bilinear extension. Antisymmetry and the Jacobi identity follow from those on $\mathfrak{g}$.
\end{proof}

For real matrix Lie algebras, one often realizes the complexification as $\mathfrak{g}+i\mathfrak{g}$ inside an ambient matrix space. This identification requires $\mathfrak{g}\cap i\mathfrak{g}=\{0\}$, as explained below.

For the classical Lie algebras, this gives the following identifications:
\begin{proposition}{Complexification of Lie Algebras}{Complexification of Lie Algebras}
  \begin{itemize}
    \item $\mathfrak{gl}(n;\mathbb{R})$ is just $M_n(\mathbb{R})$, and its complexification is $M_n(\mathbb{C}) = \mathfrak{gl}(n;\mathbb{C})$.
    \item $\mathfrak{u}(n)$ is the set of all skew-Hermitian matrices, and multiplication by $i$ gives the set of all Hermitian matrices, so its complexification is $M_n(\mathbb{C})=\mathfrak{gl}(n;\mathbb{C})$.
    \item $\mathfrak{su}(n)$ is the set of all skew-Hermitian matrices with trace zero, and multiplication by $i$ gives the set of all Hermitian matrices with trace zero, so its complexification is $sl(n;\mathbb{C})$, which is the set of all complex matrices with trace zero.
    \item $\mathfrak{sl}(n;\mathbb{R})$ is the set of all real matrices with trace zero, and its complexification is $sl(n;\mathbb{C})$, which is the set of all complex matrices with trace zero.
    \item $\mathfrak{so}(n)$ is the set of all skew-symmetric matrices, and its complexification is $so(n;\mathbb{C})$, which is the set of all complex skew-symmetric matrices.
    \item $\mathfrak{sp}(n;\mathbb{R})$ is the set of all real matrices $X$ such that $X^T J + J X = 0$, and its complexification is $\mathfrak{sp}(n;\mathbb{C})$, which is the set of all complex matrices $X$ such that $X^T J + J X = 0$.
    \item $\mathfrak{sp}(n)$ is the set of all skew-Hermitian complex matrices $X$ such that $X^T J + J X = 0$, and its complexification is $\mathfrak{sp}(n;\mathbb{C})$, which is the set of all complex matrices $X$ such that $X^T J + J X = 0$.
  \end{itemize}
\end{proposition}

\begin{remark}
  There is a subtle distinction between the abstract complexification of a real vector space and its realization inside an ambient complex vector space.

  Suppose $V$ is already a real vector subspace of a complex vector space $W$. In the definition of $V_{\mathbb{C}}$, the imaginary unit $i_{\mathrm{f}}$ is formal and is not automatically identified with the complex scalar multiplication $i$ in $W$. Indeed, making this identification may introduce additional relations between elements that are distinct in the abstract complexification.

  For example, consider $V=\mathbb{C}$ regarded as a two-dimensional real vector space, with $W=\mathbb{C}$. Its abstract complexification has complex dimension two, whereas $W$ has complex dimension one. In particular, the formal expression $i_{\mathrm{f}}\cdot 1-i$, where $i_{\mathrm{f}}$ is a formal scalar and $i$ is an element of $V$, is nonzero in $V_{\mathbb{C}}$, but becomes zero if both imaginary units are identified.

  To resolve this ambiguity, we distinguish the abstract complexification
  \begin{equation}
    V_{\mathbb{C}}=V\otimes_{\mathbb{R}}\mathbb{C}
  \end{equation}
  from the concrete complex span of $V$ in $W$. There is a natural complex-linear map
  \begin{equation}
    \begin{aligned}
      \Phi:V_{\mathbb{C}}&\longrightarrow W,\\
      v_1+i_{\mathrm{f}}v_2&\longmapsto v_1+iv_2,
    \end{aligned}
  \end{equation}
  where $i_{\mathrm{f}}$ on the left is formal, while $i$ on the right is the ordinary complex scalar multiplication in $W$.

  This map is injective if and only if
  \begin{equation}
    V\cap iV=\{0\}.
  \end{equation}
  Under this condition, we may identify the abstract complexification with its image:
  \begin{equation}
    V_{\mathbb{C}}\cong V\oplus iV\subseteq W.
  \end{equation}
  Thus, when complexifying a real Lie algebra represented by complex matrices, we must verify this condition before identifying its abstract complexification with a concrete space of complex matrices.

  For example, $\mathfrak{u}(n)$ consists of skew-Hermitian matrices, while $i\mathfrak{u}(n)$ consists of Hermitian matrices. Since a matrix that is both Hermitian and skew-Hermitian must be zero, we have
  \begin{equation}
    \mathfrak{u}(n)\cap i\mathfrak{u}(n)=\{0\}.
  \end{equation}
  Moreover, every complex matrix is uniquely expressible as the sum of a skew-Hermitian matrix and a Hermitian matrix. Consequently,
  \begin{equation}
    \mathfrak{u}(n)_{\mathbb{C}}
    \cong\mathfrak{u}(n)\oplus i\mathfrak{u}(n)
    =\mathfrak{gl}(n;\mathbb{C}).
  \end{equation}
\end{remark}

\begin{proposition}{Concrete Realization of Complexification}{Concrete Realization of Complexification}
  Suppose that $\mathfrak{g}\subseteq M_n(\mathbb{C})$ is a real Lie algebra and that
  \begin{equation}
    \mathfrak{g}\cap i\mathfrak{g}=\{0\}.
  \end{equation}
  Then the abstract complexification $\mathfrak{g}_{\mathbb{C}}$ is isomorphic, as a complex Lie algebra, to the space of complex matrices
  \begin{equation}
    \mathfrak{g}+i\mathfrak{g}
    =\{X+iY\mid X,Y\in\mathfrak{g}\}
    \subseteq M_n(\mathbb{C}).
  \end{equation}
\end{proposition}

\begin{proof}
  Consider the natural map
  \begin{equation}
    \begin{aligned}
      \Phi:\mathfrak{g}_{\mathbb{C}}
      &\longrightarrow M_n(\mathbb{C}),\\
      X+i_{\mathrm{f}}Y
      &\longmapsto X+iY,
    \end{aligned}
  \end{equation}
  where $i_{\mathrm{f}}$ denotes the formal imaginary unit in the abstract complexification, while $i$ denotes the ordinary imaginary unit acting on complex matrices.

  By the definition of the complexified bracket, $\Phi$ is a complex Lie algebra homomorphism.

  To show injectivity, suppose that
  \begin{equation}
    \Phi(X+i_{\mathrm{f}}Y)=X+iY=0.
  \end{equation}
  Then $X=-iY$, so
  \begin{equation}
    X\in\mathfrak{g}\cap i\mathfrak{g}=\{0\}.
  \end{equation}
  Hence $X=0$, and consequently $Y=0$. Thus $\ker\Phi=\{0\}$, and $\Phi$ is injective.

  Since the image of $\Phi$ is precisely $\mathfrak{g}+i\mathfrak{g}$, we conclude that
  \begin{equation}
    \boxed{
      \mathfrak{g}_{\mathbb{C}}
      \cong\mathfrak{g}\oplus i\mathfrak{g}
      \subseteq M_n(\mathbb{C}).
    }
  \end{equation}
\end{proof}

\begin{proposition}{Complexification of Homomorphisms}{Complexification of Homomorphisms}
  Let $\mathfrak{g}$ be a real Lie algebra, and $\mathfrak{h}$ a complex Lie algebra. If $\phi:\mathfrak{g}\to\mathfrak{h}$ is a real Lie algebra homomorphism, then it extends uniquely to a complex Lie algebra homomorphism $\phi_{\mathbb{C}}:\mathfrak{g}_{\mathbb{C}}\to\mathfrak{h}$, defined by
  \begin{equation}
    \phi_{\mathbb{C}}(X+iY)=\phi(X)+i\phi(Y)
  \end{equation}
\end{proposition}

\section{The Exponential Map}
\begin{definition}{Exponential Map of a Lie Group}{Exponential Map of a Lie Group}
  Let $G$ be a matrix Lie group with Lie algebra $\mathfrak{g}$. The \textbf{exponential map} of $G$ is the map
  \begin{equation}
    \exp:\mathfrak{g}\to G,\quad X\mapsto e^X.
  \end{equation}
\end{definition}

Every invertible complex matrix has a complex matrix logarithm, but this logarithm need not lie in the Lie algebra of a specified group. For example, the reflection $\operatorname{diag}(1,-1)\in O(2)$ is not the exponential of a real skew-symmetric matrix. The following theorem guarantees that elements of $G$ sufficiently close to the identity do have logarithms in $\mathfrak{g}$.

\begin{theorem}{}{}
  For $0<\epsilon<\log 2$, let $U_\epsilon=\{X\in M_n(\mathbb{C}) : \|X\|<\epsilon\}$, and let $V_\epsilon=\exp(U_\epsilon)$. Then suppose $G \subseteq GL(n;\mathbb{C})$ is a matrix Lie group with Lie algebra $\mathfrak{g}$. Then there exists $0<\epsilon<\log 2$ such that for all $A\in V_\epsilon$, $A \in G \Leftrightarrow \log A \in \mathfrak{g}$.
\end{theorem}

The bound $\epsilon<\log 2$ keeps us in a neighborhood where the exponential and the chosen matrix logarithm are inverse maps. The implication $\log A\in\mathfrak{g}\Rightarrow A\in G$ follows from the definition of $\mathfrak{g}$; we prove the converse using the following lemma.

\begin{lemma}{}{}
  Suppose $B_m \in G$ and $B_m \to I$ as $m \to \infty$. Let $Y_m = \log B_m$ which is defined for sufficiently large $m$. Suppose $Y_m \neq 0$ and $Y_m / \|Y_m\| \to Y \in M_n(\mathbb{C})$ as $m \to \infty$. Then $Y \in \mathfrak{g}$.
\end{lemma}
\begin{proof}
  For any $t \in \mathbb{R}$, we have $t Y_m / \|Y_m\| \to t Y$ as $m \to \infty$. Since $\|Y_m\| \to 0$, we find integers $k_m$ such that $k_m \|Y_m\| \to t$ as $m \to \infty$. Then we have
  \begin{equation}
    e^{k_m Y_m} = \exp\left(k_m \|Y_m\| \frac{Y_m}{\|Y_m\|}\right) \to e^{t Y}.
  \end{equation}
  As the left-hand side is in $G$ for all $m$, the right-hand side must also be in $G$. Hence $e^{t Y} \in G$ for all $t \in \mathbb{R}$, which implies that $Y \in \mathfrak{g}$.
\end{proof}

\begin{proof}[Proof of the theorem]
To prove the theorem, we separate the Lie algebra directions from their orthogonal complement. Identify $M_n(\mathbb{C}) \cong \mathbb{R}^{2n^2}$, and let $D = \mathfrak{g}^\perp$ under the Euclidean inner product. Define a map $\Phi: M_n(\mathbb{C}) \to M_n(\mathbb{C})$ by
\begin{equation}
  \Phi(Z) = e^X e^Y, \quad Z = X + Y, X \in \mathfrak{g}, Y \in D.
\end{equation}
This map is continuously differentiable, satisfies $\Phi(0)=I$, and has derivative $D\Phi_0=\operatorname{id}$. By the inverse function theorem, there exists a neighborhood $U$ of $0$ such that $\Phi$ is a diffeomorphism from $U$ onto an open set $V = \Phi(U)$ containing $I$. Choose $\epsilon$ sufficiently small that $V_\epsilon\subseteq V$.

Therefore, for any $A \in V_\epsilon$, we can write $A = e^X e^Y$ for some $X \in \mathfrak{g}$ and $Y \in D$. If $A \in G$, then we have $e^{-X} A = e^Y \in G$. As $A$ approaches $I$, we have $Y \to 0$.

Suppose, for a contradiction, that there is a sequence $Y_m\in D\setminus\{0\}$ tending to $0$ with $B_m=e^{Y_m}\in G$. Then $B_m\to I$ and, for sufficiently large $m$, $\log B_m=Y_m$. By compactness of the unit sphere, we can pass to a subsequence such that $Y_m / \|Y_m\| \to Y \in M_n(\mathbb{C})$. By the lemma, we have $Y \in \mathfrak{g}$. But $Y\in D$ and $\|Y\|=1$, contradicting $\mathfrak{g}\cap D=\{0\}$. Therefore, for sufficiently small $\epsilon$, we must have $Y = 0$, and hence $A = e^X$ with $X \in \mathfrak{g}$.
\end{proof}

\begin{corollary}{Local Homomorphism of Exponential}{Local Homomorphism of Exponential}
  Let $G$ be a matrix Lie group with Lie algebra $\mathfrak{g}$. Then there exist neighborhoods $U$ of $0$ in $\mathfrak{g}$ and $V$ of $I$ in $G$ such that $\exp:U\to V$ is a diffeomorphism, with inverse given by the matrix logarithm.
\end{corollary}


Figure~\ref{fig:exponential-local-coordinates} illustrates the corollary: the exponential map gives coordinates on a neighborhood of $I$, even when it is not globally one-to-one or onto.

\begin{figure}[H]
    \centering
    \incfig{exponential-local-coordinates}
    \caption{Near $0$ and $I$, exponentiation and the chosen logarithm are inverse coordinate maps. The highlighted arc is an open neighborhood relative to the group. The curve is schematic; the group need not be one-dimensional.}
    \label{fig:exponential-local-coordinates}
\end{figure}

\begin{remark}
  The essential difficulty in proving that the exponential map is a local diffeomorphism is not the smoothness or local invertibility of the matrix exponential itself, but rather showing that its image contains an open neighborhood of the identity in $G$.

  Indeed, the matrix exponential is already a local diffeomorphism from a neighborhood of $0$ in $M_n(\mathbb{C})$ onto a neighborhood of $I$ in $GL(n;\mathbb{C})$. Restricting this map to a sufficiently small neighborhood $U\subseteq\mathfrak{g}$ gives a diffeomorphism onto its image, and by the definition of the Lie algebra,
  \begin{equation}
    \exp(U)\subseteq G.
  \end{equation}
  However, this inclusion alone does not guarantee that $\exp(U)$ is open in $G$. In general, restricting a diffeomorphism to a subspace gives a diffeomorphism onto its image, but that image need not be open in a larger subspace containing it. For example, an open interval in a line is not open in the plane.

  The nontrivial content of the preceding theorem is precisely that, for sufficiently small $\epsilon>0$,
  \begin{equation}
    \boxed{
      G\cap V_\epsilon
      =\exp(\mathfrak{g}\cap U_\epsilon),
    }
  \end{equation}
  where $U_\epsilon$ is a neighborhood of $0$ in the ambient matrix space and $V_\epsilon=\exp(U_\epsilon)$. Thus, the exponential image is not merely contained in $G$, but actually contains a neighborhood of $I$ relative to $G$.

  Geometrically, the proof establishes this by decomposing the ambient matrix space as
  \begin{equation}
    M_n(\mathbb{C})=\mathfrak{g}\oplus D,
    \qquad D=\mathfrak{g}^{\perp},
  \end{equation}
  and showing that no nonzero sufficiently small $Y\in D$ can satisfy $e^Y\in G$. Hence there are no additional local directions of $G$ beyond those described by $\mathfrak{g}$.

  Once $G$ is equipped with the resulting smooth manifold structure, this means that $\exp(U)$ has codimension zero in $G$, or equivalently, that it is an open neighborhood of the identity.

  In the abstract Lie group approach, where $G$ is already assumed to be a smooth manifold and $\mathfrak{g}=T_I G$, the conclusion follows immediately from the inverse function theorem, since
  \begin{equation}
    (d\exp)_0=\operatorname{id}_{\mathfrak{g}}.
  \end{equation}
  In contrast, for matrix Lie groups defined merely as closed subgroups of $GL(n;\mathbb{C})$, the preceding theorem can be used to establish the smooth manifold structure itself. Thus, its significance lies in proving that $G$ is locally modeled by its Lie algebra, rather than assuming this fact in advance.
\end{remark}

\begin{corollary}{Embedded Submanifold}{Embedded Submanifold}
  Let $G$ be a matrix Lie group with Lie algebra $\mathfrak{g}$. Let $k$ be the dimension of $\mathfrak{g}$ as a real vector space. Then $G$ is a smoothly embedded submanifold of $M_n(\mathbb{C})$ of real dimension $k$, and hence a Lie group.
\end{corollary}
It follows from the corollary that $G$ is locally path connected. Therefore, $G$ is path connected if and only if it is connected.

\begin{corollary}{Tangent Space at Identity}{Tangent Space at Identity}
  Suppose $G$ is a matrix Lie group with Lie algebra $\mathfrak{g}$. Then $\mathfrak{g}$ is the tangent space of $G$ at the identity.
\end{corollary}

\begin{corollary}{Exponential Decomposition}{Exponential Decomposition}
  If $G$ is a connected matrix Lie group, every element $A \in G$ can be expressed as a finite product of exponentials of elements of $\mathfrak{g}$:
  \begin{equation}
    A = e^{X_1} e^{X_2} \cdots e^{X_m}, \quad X_i \in \mathfrak{g}.
  \end{equation}
\end{corollary}
\begin{proof}
  Let $H$ be the set of all finite products of exponentials of elements of $\mathfrak{g}$:
  \begin{equation*}
    H=\{e^{X_1}\cdots e^{X_m}:m\geq1,\ X_i\in\mathfrak{g}\}.
  \end{equation*}
  This is a subgroup of $G$. By the local diffeomorphism property of the exponential map, there is an open neighborhood $V$ of $I$ in $G$ such that every element of $V$ is an exponential of an element of $\mathfrak{g}$. Hence $V\subseteq H$. For each $h\in H$, left translation is a homeomorphism, so $hV$ is an open neighborhood of $h$ contained in $H$. Therefore
  \begin{equation*}
    H=\bigcup_{h\in H}hV
  \end{equation*}
  is open in $G$.

  Moreover, $H$ is closed. To see this, each left coset $gH$ is open, being a translate of $H$. If $g\notin H$, then $gH\cap H=\emptyset$: otherwise $gh\in H$ for some $h\in H$, which would imply $g=(gh)h^{-1}\in H$. Consequently,
  \begin{equation*}
    G\setminus H=\bigcup_{g\notin H}gH
  \end{equation*}
  is open.

  Thus $H$ is a nonempty subset of $G$ that is both open and closed. Since $G$ is connected, we must have $H=G$. By the definition of $H$, every $A\in G$ is a finite product of exponentials of elements of $\mathfrak{g}$.
\end{proof}

\begin{remark}
  More generally, a subgroup containing a nonempty open subset of a Lie group is open and closed, so it is a union of connected components containing the identity component. If the ambient group is connected, the subgroup must be the whole group.
\end{remark}

\begin{proof}[Alternative proof by subdividing a path]
  Since $G$ is connected and locally path connected, choose a continuous path $A(t)$ from $I$ to $A$. Let $V$ be an exponential neighborhood of $I$. By continuity and compactness of $[0,1]$, a sufficiently fine partition $t_j=j/m$ makes every increment
  \begin{equation*}
    A(t_{j-1})^{-1}A(t_j)\in V.
  \end{equation*}
  Write this increment as $e^{X_j}$, with $X_j\in\mathfrak{g}$. The intermediate terms cancel, giving
  \begin{equation*}
    e^{X_1}\cdots e^{X_m}
    =\bigl(A(0)^{-1}A(t_1)\bigr)\cdots
      \bigl(A(t_{m-1})^{-1}A(1)\bigr)
    =A.
  \end{equation*}
\end{proof}

\begin{remark}
  Both proofs use the fact that elements near $I$ are exponentials. The subgroup proof obtains existence from connectedness; the path proof gives the factors explicitly as $X_j=\log(A(t_{j-1})^{-1}A(t_j))$ once a path and a sufficiently fine partition are chosen. The small path segments themselves need not be exponential curves: it is their endpoint increments that are exponentials.
\end{remark}

\begin{figure}[ht]
    \centering
    \incfig{exponential-path-steps}
    \caption{Subdivide a path into small steps. Translating the endpoints of each step by $A_{j-1}^{-1}$ gives $I$ and an increment $e^{X_j}$ in the exponential neighborhood $V$. Multiplying the increments recovers $A$. The path is schematic.}
    \label{fig:exponential-path-steps}
\end{figure}

\begin{remark}
  Even if $G$ is connected, it is in general not true that every element of $G$ is the exponential of a single element of $\mathfrak{g}$.
\end{remark}

For a geometric illustration, consider $SL(2,\mathbb{R})$, which is connected. Set
\begin{equation*}
  J=\begin{pmatrix}0&-1\\1&0\end{pmatrix},\qquad
  K=\begin{pmatrix}0&1\\1&0\end{pmatrix}.
\end{equation*}
The plane $X=yJ+zK$ is a slice of $\mathfrak{sl}(2,\mathbb{R})$ through $0$. Its exponential curves lie on the surface
\begin{equation*}
  A=xI+yJ+zK,\qquad \det A=x^2+y^2-z^2=1,
\end{equation*}
a two-dimensional slice of the three-dimensional group. Each straight ray $tX$ maps to a one-parameter subgroup $A(t)=e^{tX}$ starting at $I$. Its velocity is the left translate of $X$:
\begin{equation*}
  A'(t)=A(t)X=(L_{A(t)})_*X.
\end{equation*}
Thus the curves in Figure~\ref{fig:sl2-exponential-fibers} show exponential images of rays; the arrows along the rotation curve show transported tangent vectors. These curves need not be disjoint and are not fibers in the sense of inverse images under $\exp$.

\begin{figure}[ht]
    \centering
    \incfig{sl2-exponential-fibers}
    \caption{Straight rays through $0$ become exponential curves through $I$. The red region of the hyperboloid, including its boundary except for $-I$, is unreachable by a single exponential. The drawing shows a two-dimensional slice, not the whole group.}
    \label{fig:sl2-exponential-fibers}
\end{figure}

Indeed, for any $X\in\mathfrak{sl}(2,\mathbb{R})$, the Cayley--Hamilton theorem gives $X^2=-\det(X)I$. Consequently $\operatorname{tr}(e^X)$ is either $2\cosh r$, $2$, or $2\cos r$ (by separating the even and odd terms of the exponential series), and hence is at least $-2$. Equality occurs only when $e^X=-I$. In the pictured slice, the forbidden set is therefore $x\leq-1$ with $-I$ removed; every point with $x>-1$, as well as $-I$, is an exponential.

Nevertheless, forbidden points remain accessible by products of exponentials. For example, for $s\neq0$,
\begin{equation*}
  -e^{sK}=e^{\pi J}e^{sK},\qquad \operatorname{tr}(-e^{sK})=-2\cosh s<-2.
\end{equation*}
This gives a point in the red region as a product of two exponentials, illustrating the preceding corollary.

\begin{corollary}{}{}
  Let $G$ and $H$ be matrix Lie groups with Lie algebras $\mathfrak{g}$ and $\mathfrak{h}$, respectively. Assume $G$ is connected. Suppose $\Phi_1$ and $\Phi_2$ are Lie group homomorphisms from $G$ to $H$ and that $\phi_1$ and $\phi_2$ are the corresponding Lie algebra homomorphisms from $\mathfrak{g}$ to $\mathfrak{h}$. If $\phi_1=\phi_2$, then $\Phi_1=\Phi_2$.
\end{corollary}

\end{document}
